%%% aufsteller.tex — dsa5aufsteller.sty vorfuehren: ein einfacher Bogen nur
%%% mit Groesse M, dann ein gemischter Bogen mit S, M, L und XL zusammen.
%%%
%%% Beide Bogen bekommen automatisch eine Rueckseite, deren Karten an
%%% derselben physischen Blattstelle stehen wie vorne -- siehe die
%%% Begruendung bei \dsaAufstellerbogen in ../dsa5aufsteller.sty.
%%%
%%% Platzhalterbild fuer alle Karten: grafiken/titelbild.jpg, die einzige
%%% Grafik im Projekt, die nicht aus dem Ulisses-Baukasten stammt.
%%%
%%%   xelatex aufsteller.tex
%%%   python3 ../werkzeuge/nachmessen.py aufsteller.pdf

\documentclass[ohnehintergrund]{dsa5latex}
\graphicspath{{../}}
\usepackage{dsa5aufsteller}

\begin{document}

% Bogen 1: nur Groesse M, sauberes 3x2-Raster.
\begin{dsaAufstellerbogen}
\dsaAufstellerM{0mm}{0mm}{\dsagrafik{titelbild}}{Ork}
\dsaAufstellerM{40mm}{0mm}{\dsagrafik{titelbild}}{Goblin}
\dsaAufstellerM{80mm}{0mm}{\dsagrafik{titelbild}}{Wache}
\dsaAufstellerM{0mm}{60mm}{\dsagrafik{titelbild}}{Baeuerin}
\dsaAufstellerM{40mm}{60mm}{\dsagrafik{titelbild}}{Bandit}
\dsaAufstellerM{80mm}{60mm}{\dsagrafik{titelbild}}{Kultistin}
\end{dsaAufstellerbogen}

% Bogen 2: gemischte Groessenklassen auf einem Bogen, frei platziert --
% genau der Fall, den ein starres Spaltenraster nicht koennte. Alle vier
% Klassen stehen hochkant (Bogen oben, gerade Kante unten wie bei M).
\begin{dsaAufstellerbogen}
\dsaAufstellerXL{0mm}{0mm}{\dsagrafik{titelbild}}{Drache}
\dsaAufstellerL{85mm}{0mm}{\dsagrafik{titelbild}}{Troll}
\dsaAufstellerM{85mm}{70mm}{\dsagrafik{titelbild}}{Waldlaeufer}
\dsaAufstellerS{0mm}{106mm}{\dsagrafik{titelbild}}{Rabe}
\dsaAufstellerS{26mm}{106mm}{\dsagrafik{titelbild}}{Fuchs}
\dsaAufstellerS{52mm}{106mm}{\dsagrafik{titelbild}}{Kroete}
\end{dsaAufstellerbogen}

\end{document}
